% ===== 原文 pp.150–152（PDF p.160–162） =====
% [原文 p.150]

\chapter{第 9 章中命题的证明}

在本附录中，我们给出第 9 章正文中所讨论的部分选定命题的证明。

\begin{proposition}[1] 若 $(Y,F)$ 满足：对所有允许对 $(u,v)$ 都有 $G(u,u)=-G(v,v)$，则 $v\in TY$，且 $(v,u)$ 也是允许对。\end{proposition}

\begin{proof}
首先我们证明 $Y$ 必定是余迷向子流形（采用 §9.3 的术语）。这是因为 $v_0\in(TY)^\perp$ 蕴含 $(0,v_0)$ 是允许对。于是
\[
  G(v_0,v_0)=-G(0,0)=0 .
\]
这又蕴含 $(TY)^\perp$ 是迷向子空间，因为
\[
  G(v_0,v_0')=\frac{1}{2}G(v_0+v_0',v_0+v_0')=0
\]
对所有 $v_0,v_0'\in(TY)^\perp$ 成立。也就是说，$(TY)^\perp\subseteq(TY^\perp)^\perp=TY$，而这正是 $Y$ 为余迷向的含义。

% [Tutor Note: p.150 两处省略的论证补出。(i) 为何 $v_0\in(TY)^\perp$ 蕴含 $(0,v_0)$ 是允许对：按第 9 章定义 1，允许对的第一分量必须落在 $TY$ 中，而 $0\in TY$；把允许对 $(u,v)=(0,v_0)$ 代入假设 $G(u,u)=-G(v,v)$ 即得 $G(v_0,v_0)=-G(0,0)=0$。(ii) 中间的极化恒等式由双线性展开得到：$G(v_0+v_0',v_0+v_0')=G(v_0,v_0)+2G(v_0,v_0')+G(v_0',v_0')$，而 $G(v_0,v_0)=G(v_0',v_0')=0$；右端为零则是因为 $(0,v_0+v_0')$ 也是允许对（$v_0+v_0'\in(TY)^\perp$），再用同一假设。]

其次我们证明：若 $(u,v)$ 是允许对，则 $v$ 落在子空间 $TY$ 中。（允许对的定义要求 $u\in TY$，但先验地并不要求 $v\in TY$。）设 $v_0\in(TY)^\perp$。对实参数 $t$ 的任意取值，容易看出 $(u,v+tv_0)$ 也是允许对。因此
\[
  -G(u,u)=G(v+tv_0,v+tv_0)
\]
对所有 $t$ 成立。对上式两边关于 $t$ 求导，得到
\[
  0=\frac{\dd}{\dd t}G(v+tv_0,v+tv_0)=2G(v,v_0)
\]
（因为 $G(v,v)$ 与 $t$ 无关，且 $G(v_0,v_0)=0$）。$v_0\in(TY)^\perp$ 的选取是任意的，所以
\[
  v\in((TY)^\perp)^\perp=TY .
\]

% [Tutor Note: p.150 “容易看出 $(u,v+tv_0)$ 也是允许对”的省略论证：允许对条件（第 9 章定义 1）是关于一切 $w\in TY$ 的一组线性等式；由于 $v_0\in(TY)^\perp$，$v_0$ 与 $TY$ 的配对按 $G$ 的正交补定义为零，且由 $(0,v_0)$ 是允许对又知相应的 $G(v_0,w)$ 也为零（$v_0\in(TY)^\perp$ 的 G-正交性），故把 $v$ 换成 $v+tv_0$ 不改变这组等式，$(u,v+tv_0)$ 仍是允许对。]

% [原文 p.151]
最后我们证明：$(u,v)$ 是允许对蕴含 $(v,u)$ 也是允许对。由线性性，若 $(u,v)$ 与 $(u',v')$ 是允许对，则 $(u+u',v+v')$ 也是允许对。按假设，
\[
  G(u+u',u+u')=-G(v+v',v+v')
\]
利用 $G(u,u)=-G(u',u')$ 与 $G(u',u')=-G(v',v')$，这蕴含
\begin{equation}
  G(u,u')=-G(v,v') . \tag{D.1}
\end{equation}
设 $(u,v)$ 是允许对，$w\in TY$ 任意。不难看出存在 $v'\in TY$ 使得 $(w,v')$ 是允许对。在式 (D.1) 中取 $u'=w$，可得
\[
  G(u,w)=-G(v,v')=-G(v',v)=-F(w,v)=F(v,w).
\]
我们由此断言 $(v,u)$ 也是允许对。
\end{proof}

% [Tutor Note: p.151 式 (D.1) 的推导与原文笔误。正文照原文排印：原文写 “using $G(u,u)=-G(u',u')$ and $G(u',u')=-G(v',v')$”（PDF 图像与文字层一致，并非抽取错误）。但代入 $G(u+u',u+u')=G(u,u)+2G(u,u')+G(u',u')$ 与 $G(v+v',v+v')=G(v,v)+2G(v,v')+G(v',v')$ 后，要消去对角项必须用 $G(u,u)=-G(v,v)$（而不是 $-G(u',u')$），才得到 $2G(u,u')=-2G(v,v')$，即 (D.1)。故原文该处应为 $G(u,u)=-G(v,v)$；本译本照原文排印，正确形式在此批注给出。]

% [Tutor Note: p.151 的 “Theorem 1” 是原文的悬空引用。全文检索确认本书没有编号为「定理 1」的定理（第 9 章只有定义 1–7 与命题 1–10；附录 C 的定理均不编号）。按上下文，此处最可能想指第 9 章中给出允许对对称性的命题 1，或本附录的命题 (1)。译者照原文直译为「定理 1」，未作擅改。]

\begin{proposition}[3] $(Y,F)$ 定义一个拓扑双膜，当且仅当 $Y\subseteq X\times\widehat X$ 是余迷向子流形，满足 $\ker F=(TY)^\perp$（即 $F$ 的退化方向与 $G|_{TY}$ 的退化方向一致），并且此外在 $SY$ 上
\[
  (\widetilde G^{-1}\widetilde F)^2=+1,
\]
其中 $\widetilde G\equiv G|_{SY}$、$\widetilde F\equiv F|_{SY}$。\end{proposition}

\begin{proof}
为证明“仅当”，假设 $(Y,F)$ 是拓扑双膜。定理 1 告诉我们：若 $(u,0)$ 是允许对，则 $(0,u)$ 也是允许对，反之亦然；这蕴含 $\ker F=(TY)^\perp$。于是 $F$ 的退化方向与 $G|_{TY}$ 的退化方向一致，且 $\widetilde F$ 与 $\widetilde G$（$G$ 与 $F$ 在屏幕分布 $SY$ 上的限制）都是非退化的。

此外，若 $\tilde u,\tilde v\in SY$，则按定义，$(\tilde u,\tilde v)$ 是允许对当且仅当 $\tilde v=(\widetilde G^{-1}\widetilde F)\tilde u$。由于 $(u,(\widetilde G^{-1}\widetilde F)\tilde u)$ 是允许对，$((\widetilde G^{-1}\widetilde F)\tilde u,u)$ 也是允许对（再次由定理 1）。这意味着 $\tilde u=(\widetilde G^{-1}\widetilde F)^2\tilde u$；由于 $\tilde u\in SY$ 任意，我们有
\[
  (\widetilde G^{-1}\widetilde F)^2=1 .
\]

% [Tutor Note: p.151 原文此处写 $(u,(\widetilde G^{-1}\widetilde F)\tilde u)$ 是允许对，其中第一个分量印作 $u$ 而非 $\tilde u$，疑为笔误。依据：(i) 紧接的上一句刚说明 $(\tilde u,\tilde v)$ 是允许对当且仅当 $\tilde v=(\widetilde G^{-1}\widetilde F)\tilde u$，允许对的第一个分量应取 $\tilde u\in SY$；(ii) 结论 $\tilde u=(\widetilde G^{-1}\widetilde F)^2\tilde u$ 只能由 $(\tilde u,(\widetilde G^{-1}\widetilde F)\tilde u)$ 经定理 1 交换次序、再用一次上述定义得到。译文照印文保留 $u$，读者应理解为 $\tilde u$。]

为证明“若”，设 $(u,v)$ 是允许对；我们想证明 $v\in TY$ 且 $(v,u)$ 也是允许对。设 $v_0\in(TY)^\perp$ 任意。我们于是有
\[
  G(v,v_0)=F(v,v_0)=0
\]

% [原文 p.152]
其中第一个等号由允许对的定义给出，第二个等号由 $\ker F=(TY)^\perp$ 这一事实给出。因此 $v\in((TY)^\perp)^\perp=TY$。

% [Tutor Note: p.161 末行（续至 p.152 首行）的等式按印文为 $G(v,v_0)=F(v,v_0)=0$，疑为笔误，应为 $G(v,v_0)=F(u,v_0)=0$。依据：(i) 本证明稍后明确以 $F(u,w)=G(v,w)$ 作为允许对条件（p.152），故取 $w=v_0$ 应得 $F(u,v_0)$；(ii) 原文对两个等号的解释（第一个用允许对的定义、第二个用 $\ker F=(TY)^\perp$）只有当第二项的第一分量为 $u\in TY$ 时才成立，因为 $\ker F=(TY)^\perp$ 只保证 $F$ 在 $TY$ 与 $(TY)^\perp$ 之间的配对上为零；(iii) 要得到结论 $v\in((TY)^\perp)^\perp=TY$，必须先用 $G(v,v_0)=0$ 判定 $v\in TY$。译文照印文保留。]

现在我们把 $TY$ 中的向量按分解 $TY=(TY)^\perp\oplus SY$ 写成 $u=u_0+\tilde u$ 与 $v=v_0+\tilde v$。由于 $F(u,w)=G(v,w)$ 对一切 $w\in TY$ 成立，我们有 $F(\tilde u,\tilde w)=G(\tilde v,\tilde w)$ 对一切 $\tilde w\in SY$ 成立，即 $\tilde v=(\widetilde G^{-1}\widetilde F)\tilde u$。但 $(\widetilde G^{-1}\widetilde F)^2=1$，故 $\tilde u=(\widetilde G^{-1}\widetilde F)\tilde v$。于是 $F(\tilde v,\tilde w)=G(\tilde u,\tilde w)$ 对一切 $\tilde w\in SY$ 成立，并且最终 $F(v,w)=G(u,w)$ 对一切 $w\in TY$ 成立。\end{proof}

\begin{proposition}[6] 子丛 $\cR_+$（如上面所定义）关于 $e^*TM$ 是拉格朗日的。特别地，它是 $n$ 维的。\end{proposition}

\begin{proof}
先细化命题 3 中给出的拓扑双膜切丛的描述会比较方便。设 $\widetilde R$ 为自同态 $\widetilde G^{-1}\widetilde F:SY\to SY$，由命题 3，它平方等于恒等。于是屏幕分布可分解为 $\widetilde R$ 的 $+1$ 与 $-1$ 特征子丛：$SY=\widetilde{\cR}_+\oplus\widetilde{\cR}_-$。事实上，由于 $\widetilde G^{-1}\widetilde F$ 是一个对称矩阵与一个反对称矩阵的乘积，它是无迹的，特别地其特征值恰有一半是 $+1$、一半是 $-1$。因此
\[
  \dim\widetilde{\cR}_+=\dim\widetilde{\cR}_-=\frac{1}{2}\dim SY=\frac{1}{2}(2k-2n)=k-n,
\]
其中 $\dim M=2n$、$\dim Y=k$（回忆：由于 $Y$ 是余迷向的，$k\ge n$）。由 $\cR_+$ 的定义可以直接看出
\[
  \cR_+=(TY)^\perp\oplus\widetilde{\cR}_+ .
\]
因此其维数为
\[
  \dim\cR_+=\dim (TY)^\perp+\dim\widetilde{\cR}_+=(2n-k)+(k-n)=n .
\]
此外，若 $s,t\in\cR_+$，则
\[
  G(s,t)=F(s,t)=-F(t,s)=-G(t,s)=-G(s,t)=0,
\]
所以 $\cR_+$ 是 $e^*TM$ 的迷向子丛。（在上式的第一与第三个等号中我们用了 $\cR_+$ 的定义，在第二与第四个等号中我们用了 $F$ 与 $G$ 的对称性。）由于它是中间维数且迷向的，由此可得 $\cR_+$ 关于 $e^*TM$ 是拉格朗日的。\end{proof}

% [Tutor Note: p.152 “无迹”与维数计数补出。设 $H=\widetilde G^{-1}\widetilde F$；因 $\widetilde G$ 对称正定，$H$ 与 $\widetilde G^{-1/2}\widetilde F\widetilde G^{-1/2}$ 相似，而后者（对称矩阵乘反对称矩阵）的转置是自身的负，故无迹；平方为 $1$ 且无迹 $\Rightarrow$ $+1$ 与 $-1$ 特征值各占一半。维数核对：$\dim SY=\dim TY-\dim (TY)^\perp=k-(2n-k)=2k-2n$，故 $\dim\widetilde{\cR}_\pm=k-n$；再由 $\cR_+=(TY)^\perp\oplus\widetilde{\cR}_+$ 得 $\dim\cR_+=(2n-k)+(k-n)=n=\frac{1}{2}\dim M$，配合迷向性即得 $\cR_+$ 是拉格朗日的。]
